% HD-G27. GPT-6 Astra Pro /root/g34, inherited effort; 21 November 2026.
\section{HD-G27: reflection-containing sources over the seven remaining fields}
\label{sec:other-fields-reflection-scope}

Put
\[
 K=\mathbf Q(\sqrt{-d}),\qquad
                 d\in\{2,7,11,19,43,67,163\}.                 \tag{G27.1}
\]
Let \(G\subset GL_m(K)\) be an actual finite absolutely irreducible
matrix group whose character field is exactly \(K\). Consider its
natural primitive rational group-algebra Jacobian factor in a
connected full-\(G\) sphere-cover family, with the indicated
\(K\)-corner and restricted theta polarization. Every such Weil
sixfold of signature \((3,3)\) whose period map has rank nine,
and whose coefficient group contains a reflection, has ordinary
Hermitian determinant class \([-1]\).

More precisely, the faithful irreducible reflection block in the
proof is either \(G_{24}\) over \(\mathbf Q(\sqrt{-7})\), or the
explicit order-\(2304\) group
\[
 \left\langle W(F_4),\frac{\sqrt{-2}}2
 \begin{pmatrix}1&1&0&0\\1&-1&0&0\\0&0&1&1\\0&0&1&-1\end{pmatrix}
 \right\rangle
             \quad\hbox{over }\mathbf Q(\sqrt{-2}).           \tag{G27.2}
\]
Consequently no reflection-containing rank-nine source of this
specified kind occurs for the other five fields in (G27.1).
Every source with \(m\le8\) contains a reflection and is therefore
covered by these conclusions. The proof supplies the actual
induction, quotient-cover, primitive-corner and polarization maps.
It makes no assertion about reflection-free groups of larger rank.

\paragraph{1. Normal reflection blocks over the actual field.}
Every field in (G27.1) has \(\mu(K)=\{1,-1\}\).
For completeness, a root of unity in a quadratic field has
\(\varphi(q)\le2\), hence order \(q\in\{1,2,3,4,6\}\).
The last three nonreal possibilities generate
\(\mathbf Q(i)\) or \(\mathbf Q(\sqrt{-3})\), neither in (G27.1).

Average a positive \(K\)-Hermitian form \(h\) over \(G\).
Let \(W\) be the nontrivial normal subgroup generated by all
reflections in \(G\). Join their \(K\)-root lines
\(\operatorname{im}(s-1)\) when they are not orthogonal.
The spans of different connected components are orthogonal.
Their sum is a nonzero \(G\)-invariant subspace, hence all of \(V\).
Each component span is an absolutely irreducible reflection
module: an invariant complex subspace and its orthogonal
complement place each rank-one image wholly in one of them,
and connectedness then places the entire component there.
The group \(G\) permutes the components transitively.

Choose one component \(U\), put \(r=\dim_KU\), and let
\(H\) be its stabilizer. Write \(\sigma:H\to GL(U)\),
\(H_0=\ker\sigma\), \(G'=H/H_0\), and \([G:H]=a\).
For left coset representatives \(t_j\), the literal isomorphism
\[
 \bigoplus_{j=1}^aU\longrightarrow V,\qquad
       (u_j)\longmapsto\sum_j\rho(t_j)u_j,\qquad m=ar,        \tag{G27.3}
\]
has block matrix \(\sigma(h)\) when \(gt_j=t_kh\).
It identifies the coefficient form with the orthogonal sum of
copies of \(h_U\), without a scalar. Its induced trace formula
and the actual \(K\)-matrices give
\[
 K=\mathbf Q(\operatorname{tr}\rho)
   \subseteq\mathbf Q(\operatorname{tr}\sigma)\subseteq K.
                                                               \tag{G27.4}
\]
In particular the faithful block has exact character field \(K\).
For \(r=1\), its character takes values in \(\mu(K)=\{\pm1\}\).
Formula (G27.3) would then make every trace of \(\rho\) rational,
contradicting (G27.4). Thus \(r\ge2\).

\paragraph{2. The quotient curves, exact primitive factor and period bound.}
For the original cover \(\widetilde C\to\mathbf P^1\), use the
actual quotients
\[
             Y=\widetilde C/H,\qquad C'=\widetilde C/H_0.
                                                               \tag{G27.5}
\]
Then \(C'\to Y\) has group \(G'\). With left deck actions,
the sheets of \(Y\) are \(H\backslash G\); the sheet
\(Ht_j^{-1}\) corresponds to the block \(t_jH\) in (G27.3).
The explicit coefficient map and its inverse are
\[
 F_v(x)=\pi_U\rho(x)v,\quad F_v(hx)=\sigma(h)F_v(x),\quad
 F_{\rho(g)v}(x)=F_v(xg),\qquad
 F\longmapsto\sum_{x\in H\backslash G}\rho(x)^{-1}F(x).
                                                               \tag{G27.6}
\]
Each summand in the inverse is independent of its representative.
Taking \(x=t_j^{-1}\) proves it is the inverse of (G27.3).
A cycle of sheets at a puncture has invariant stalk equal to
the fixed space of the product of its block monodromies.
Thus (G27.6) identifies the middle extensions, including ramified
sheets. The finite-sheet cellular map gives the same first
parabolic cohomology and preserves cup product followed by
the trace of sheets.

The equality of the abelian factors is also literal.
For \(v\in U(K)\setminus\{0\}\) and \(q=h(v,v)>0\), put
\[
 e_G=\frac{2m}{|G|q}\sum_{g\in G}
        \operatorname{Re}h(v,\rho(g^{-1})v)\,g,\qquad
 e_H=\frac{2r}{|H|q}\sum_{h\in H}
        \operatorname{Re}h_U(v,\sigma(h^{-1})v)\,h .
                                                               \tag{G27.7}
\]
The coefficients outside \(H\) vanish by orthogonality, and
\(m/|G|=r/|H|\). Hence \(e_G=e_H\) in \(\mathbf Q[G]\).
Schur orthogonality gives the rank-one projector on the selected
coefficient line and its conjugate. After the same integer clears
the denominators, their image is the same abelian subvariety
of \(J\widetilde C\), with the same restricted theta form.
The two real Schur constants agree:
\[
                    \frac{2|G|}{mq}=\frac{2|H|}{rq}.          \tag{G27.8}
\]
The element \(e_H\) is the \(H_0\)-average times the corresponding
primitive idempotent for \(G'\). For the actual quotient
\(\alpha:\widetilde C\to C'\), its factor \(B'\subset JC'\)
therefore maps isogenously onto the same factor \(B\), by
\(\alpha^*\), and
\[
 \alpha_*\alpha^*=[|H_0|],\qquad
       (\alpha^*)^*\theta_{J\widetilde C}=|H_0|\theta_{JC'}.
                                                               \tag{G27.9}
\]
The corner action is identified by the coefficient and class-sum
maps in paragraph 7 below. Its rank-six determinant multiplier is
\(|H_0|^6=N_{K/\mathbf Q}(|H_0|^3)\).
These fixed Hodge and lattice maps identify the period maps after
finite base change fixing the monodromy labels. Thus the original
rank is bounded by the pointed-curve parameter dimension of
the faithful cover \(C'\to Y\).

Let \(g=g(Y)\) and let \(b'\) count its nontrivial coefficient
monodromies, of codimensions \(c'_j\). The exact parabolic Euler
formula is
\[
                    6=r(2g-2)+\sum_{j=1}^{b'}c'_j.           \tag{G27.10}
\]
Global invariants and dual invariants vanish by irreducibility.
For fixed finite monodromy the parameter space is finite over
the moduli of the \(b'\)-pointed base curve. Its dimension is
\(3g-3+b'\) in the stable cases, \(b'\) for \(g=1,b'>0\),
at most one for an unpointed elliptic base, and \(b'-3\)
for a sphere base. An irreducible rank-at-least-two sphere
system has \(b'\ge3\), since the smaller punctured-sphere
groups are cyclic or trivial.

For \(r=2\), (G27.10) gives period bounds \(7,6,5\) at
\(g=0,1,2\), and makes larger genus impossible.
For \(r=3\), genus one gives at most six parameters and
genus two gives three; rank nine therefore forces \(g=0\),
\(b'=12\), and all \(c'_j=1\).
For \(r\ge4\), genus at least two is impossible and genus one
has at most six parameters, so again \(g=0\).

\paragraph{3. The published classification restricted by actual traces.}
Cohen, \emph{Finite complex reflection groups},
Ann.\ Sci.\ ENS 9 (1976),379--436, gives the imprimitive list
in Theorem (2.4), pp.386--387, the real list in (4.4)(i),
p.399, and the nonreal primitive list in Theorem (5.12),
p.435, with table (4.16), p.412. The original source is
\url{https://www.numdam.org/item/ASENS_1976_4_9_3_379_0.pdf}.
As defined in (1.1), p.380, the word complex there means
nonreal. The following actual traces give the required restriction.

For \(G(\ell,p,r)\), \(r\ge3\), its subgroup
\(G(\ell,\ell,r)\) contains the two matrices
\[
 \operatorname{diag}(\zeta,\zeta^{-1},1,\ldots),\qquad
 \operatorname{diag}(\zeta,\zeta,\zeta^{-2},1,\ldots).
\]
The first trace forces \(\zeta+\zeta^{-1}\) to be a rational
algebraic integer in \([-2,2]\); thus
\(\ell\in\{1,2,3,4,6\}\).
The second has trace \(r-3+3\zeta\) for \(\ell=3\),
\(r-4+2i\) for \(\ell=4\), and \(r-3+\zeta\) for
\(\ell=6\). These exclude all three for (G27.1).
The case \(\ell=1\) is the reducible permutation representation.
Only the rational signed-permutation groups remain.

For the nonreal primitive list, the actual trace tests are
\[
\begin{array}{c|c}
G_{24}&(-1+\sqrt{-7})/2\\
G_{25},G_{26}&2+\omega\ \hbox{or }2+\bar\omega\\
G_{27}&(1+\sqrt5)/2\\
G_{29},G_{31}&4i\\
G_{32}&3+\omega\ \hbox{or }3+\bar\omega\\
G_{33}&2+3\omega\\
G_{34}&6\omega
\end{array}                                                    \tag{G27.11}
\]
Here \(\omega^2+\omega+1=0\).
The \(G_{24}\) word is \((s_1s_2s_3)^2\) in paragraph 5;
its explicit intertwiner to Cohen's p.407 roots is recorded
in HD-S13. The three order-three-reflection traces follow from
their one nontrivial eigenvalue in Cohen's diagrams and table.
The \(G_{27}\) labelled five-edge, p.408, gives the product
of two reflections with eigenvalues \(e^{\pm2\pi i/5}\).
For \(G_{29},G_{31}\), (1.4)(ii), p.381, and their table
degrees give scalar centre of order four, hence the element
\(iI_4\). The displayed Gaussian matrices of HD-S32 verify
that same scalar directly.
For \(G_{33}\), Cohen (5.10), p.430, contains
\(G(3,3,4)\), whose element
\(\operatorname{diag}(\omega,\omega,\omega,1)\) fixes the
remaining coordinate and has the stated trace.
For \(G_{34}\), its scalar centre has order six by the same
degree formula and contains \(\omega I_6\).
Thus the only retained nonreal block of rank at least three
is \(G_{24}\), and only over \(\mathbf Q(\sqrt{-7})\).

The five-edge also excludes the real \(H_3,H_4\):
its two-reflection product has trace
\(r-\tfrac52+\tfrac12\sqrt5\), outside every imaginary
quadratic field. The real possibilities of rank at least three
are therefore exactly
\[
                       A_r,B_r,D_r,E_6,E_7,E_8,F_4.          \tag{G27.12}
\]
All have their standard rational root models.

Every complex equivalence just used is replaced by an actual
\(K\)-intertwiner: solve the finite equations
\(T\rho(w)=\rho_0(w)T\) over \(K\) for the fixed group
identification. The complex solution contains an invertible map,
and scalar extension preserves the linear-system rank.
Any nonzero solution is invertible by irreducibility.
The pullback of the standard invariant Hermitian form is a
positive rational scalar \(a\) times the given one, by its
scalar commutant and \(K\cap\mathbf R=\mathbf Q\).
Applying \(T\) on every parabolic cochain gives that exact
paired Hodge map. Paragraph 7 retains its Schur factors and
the integer clearing its lattices.

\paragraph{4. All real normalizers and the exceptional field \(\mathbf Q(\sqrt{-2})\).}
Let \(W_U\) be a real irreducible reflection block of rank
\(r\ge4\), normal in \(G'\). In its rational root model,
\(g^{-1}\bar g\) commutes with \(W_U\). Indeed \(g\) and
\(\bar g\) induce the same conjugation on the real matrices
of \(W_U\). Its scalar commutant gives
\(g=aR\), with \(R\) real orthogonal and \(|a|=1\).
The real root-line normalizer is finite: its action on the
finite spanning line set has kernel \(\{\pm I\}\), by
nonorthogonality connectedness. Hence \(R\) has finite order,
and the finite order of \(g\) makes \(a\) a root of unity.

Here are the actual real normalizer maps. For \(A_r\),
\(r\ge4\), its root-line graph is the line graph of the
complete graph on \(r+1\) vertices. The maximal cliques
of size \(r\) are the vertex stars; triangles have size three.
Their permutation gives a literal vertex permutation in the
augmentation representation. After removing it, every root
line is fixed and connectedness makes the map scalar.
Thus the real normalizer is rational.
For \(B_r\), the coordinate reflection class has size \(r\)
and the other class has size \(r(r-1)\); their distinct sizes
preserve the coordinate lines. Orthogonality then gives a
signed permutation, again rational.
For \(D_r\), \(r\ge5\), the root-system diagram symmetry
is the literal coordinate sign change at one fork end;
its normalizer is contained in \(W(B_r)\).
The actual chamber-to-diagram map is the one of Humphreys,
\emph{Introduction to Lie Algebras and Representation Theory},
Theorem10.3, pp.51--52, and Section12.2, pp.65--66.
The \(E_7,E_8\) diagrams have no nontrivial symmetry.
For \(E_6\), in the simple basis with edges
\((1,3),(3,4),(4,5),(5,6),(2,4)\), its diagram involution
is the literal permutation \((1\ 6)(3\ 5)\).
All these Weyl and diagram matrices are rational.

For \(D_4,F_4\), use the actual coordinate sets
\[
 \mathcal S=\{\pm e_j\}\cup
       \{\tfrac12(\pm e_1\pm e_2\pm e_3\pm e_4)\},\qquad
 \mathcal L=\{\pm e_j\pm e_k:j<k\}.                         \tag{G27.13}
\]
The orthogonal four-frames in either set have successive
counts \(24,6,4,2\). These follow by the literal orthogonality
equations after choosing the first vector; all frames were
enumerated in HD-S35. The group \(W(F_4)\) of order1152
attains the resulting bound and preserves both sets.
Thus the real normalizer of \(W(D_4)\) is \(W(F_4)\),
with rational matrices.
On the twenty-four normalized \(F_4\) root lines, the
inner-product-\(1/\sqrt2\) graph is connected bipartite,
with short and long lines its two parts. This follows from
(G27.13): every vertex has six neighbours in the opposite
part and shared-coordinate edges connect the graph.
A part-preserving normalizer has at most1152 choices by
the same frames and is exactly \(W(F_4)\).
Writing \(M\) for the integer matrix in (G27.2), the exact
identities \(M\mathcal S=\mathcal L\) and
\(\frac12M\mathcal L=\mathcal S\) show that
\(u=M/\sqrt2\) exchanges the parts. Consequently
\[
           N_{O_4(\mathbf R)}(W(F_4))
                     =W(F_4)\ \cup\ uW(F_4).                \tag{G27.14}
\]

For every rational representative \(R\) above, \(aR\in GL_r(K)\)
forces \(a\in K\). Hence \(a\in\{\pm1\}\), and the whole
group has rational character, contradicting (G27.4).
Only the second coset of (G27.14) remains.
For \(g=auw\), a nonzero rational entry of \(Mw\) gives
\(a/\sqrt2\in K\). Thus \(a^2\in K\cap\mu_\infty=\{\pm1\}\).
The value \(a^2=1\) would give the real number
\(1/\sqrt2\in K\), impossible. The value \(a^2=-1\)
gives
\[
          a=\pm i,\qquad a/\sqrt2=\pm\sqrt{-2}/2,
\]
so \(K=\mathbf Q(\sqrt{-2})\) and \(g=\pm Uw\), with
\(U=(\sqrt{-2}/2)M\). Here \(U^2=-I\in W(F_4)\).
Since \(G'\) contains \(W(F_4)\), any such element gives
exactly the two cosets in (G27.2). This proves the claimed
order2304 and exhausts the real normalizer alternatives.

The group has exact field \(K\): for
\(w=\operatorname{diag}(1,-1,1,1)\in W(F_4)\),
\(\operatorname{tr}(Uw)=\sqrt{-2}\).
Its matrices lie in \(\tfrac12\mathcal O_K\);
the \(W\)-columns are in \(\mathcal S\), and
\(M\mathcal S=\mathcal L\) gives the same denominator
for the second coset.

\paragraph{5. Every twelve-reflection \(G_{24}\) source is split.}
For \(r=3\), paragraph2 makes \(G'\) generated by twelve
reflections, so it is itself a reflection group. Formula
(G27.11), the rational alternatives, and exact character
field leave precisely \(G_{24}\) over \(\mathbf Q(\sqrt{-7})\).
Put \(t=(-1+\sqrt{-7})/2\). Its actual model is generated by
\[
 s_1=\begin{pmatrix}-1&1&t\\0&1&0\\0&0&1\end{pmatrix},\quad
 s_2=\begin{pmatrix}1&0&0\\1&-1&1\\0&0&1\end{pmatrix},\quad
 s_3=\begin{pmatrix}1&0&0\\0&1&0\\-1-t&1&-1\end{pmatrix}.
                                                               \tag{G27.15}
\]
It preserves the integral Hermitian matrix
\[
 H_0=\begin{pmatrix}2&-1&-t\\-1&2&-1\\1+t&-1&2\end{pmatrix},
                 \qquad \det H_0=1.                         \tag{G27.16}
\]
Its leading minors are \(2,3,1\).
The explicit root intertwiner of HD-S13 identifies this
with Cohen's p.407 model; its exact order336 enumeration
also proves that \(-I\) is present and that every reflection
has order two. The latter also follows from its trace and
determinant over \(K\), whose only roots of unity are \(\pm1\).
The word \((s_1s_2s_3)^2\) has trace \(t\).

Let \(R=\mathcal O_K[1/14]\) and \(L=R^3\).
The lattice (G27.16) is Hermitian self-dual, every boundary
inertia has the integral projector \((1+g_j)/2\), and
\(-I\) kills global coinvariants for \(L\) and its dual.
The perfect parabolic construction of paragraph6 therefore
gives a perfect rank-six skew-Hermitian cup matrix \(S\)
over \(R\). Multiplication by \(s=\sqrt{-7}\), a unit of
\(R\), gives the Hermitian matrix \(sS\).
Its negative rational determinant is a unit of \(R\), so
\[
                  \det(sS)=-2^a7^b,\qquad a,b\in\mathbf Z.
                                                               \tag{G27.17}
\]
Both primes are norms:
\(N(t)=2\) and \(N(2t+1)=7\).
Thus every such tuple, not just the tuple used to prove
dominance in HD-G12, has ordinary determinant class \([-1]\).

\paragraph{6. The \(F_4\) profile and the perfect integral pairing.}
For (G27.2), paragraph2 makes the base a sphere.
The parabolic rank-six formula gives
\(\sum c_j=14\), and period rank nine requires \(b'\ge12\).
The quotient by \(W(F_4)\) has order two, so a generating
product-one tuple has at least two entries in its other coset.
Each such element is \(iR\) with \(R\) real orthogonal;
its fixed space is the \(-i\)-eigenspace of \(R\).
Conjugation pairs this with the equally large \(i\)-eigenspace,
so its codimension is at least two.
These two entries use all of
\(\sum(c_j-1)=14-b'\le2\).
Consequently \(b'=12\), exactly two entries have codimension
two, and the other ten are real reflections.
For either exceptional entry its two eigenspaces exhaust the
four-dimensional space, hence \(R^2=-I\) and \((iR)^2=I\).
All twelve local inertias are therefore involutions.

Now take \(R=\mathcal O_{\mathbf Q(\sqrt{-2})}[1/2]\),
\(L=R^4\), and coefficient Gram matrix \(I_4\).
The actual half-integral matrices preserve this self-dual
lattice. The global \(-I\) and the same local involution
projectors give the following complete integral argument.

Let \(\Sigma\) be the sphere with branch discs removed.
Its first homology is the kernel of a map of finite free
\(R\)-modules because \(\Sigma\) retracts onto a graph.
For \(L\) and its dual, \(H_0(\Sigma,L)=0\): the word
\(-I\) gives the unit \(2\) in the sum of the monodromy
differences. Universal coefficients and the actual cap-product
Poincare--Lefschetz maps give perfect pairings of absolute
and relative first homology with the dual coefficients.
These modules are free. The rings here and in paragraph5
are principal ideal domains: \(\mathbf Z[\sqrt{-2}]\) is
norm-Euclidean by nearest-coordinate bound \(3/4\);
for \(\mathbf Z[t]\), write a quotient as \(x+yt\),
choose \(b\in\mathbf Z\) nearest \(y\), and then
\(a\in\mathbf Z\) nearest \(x-(y-b)/2\).
The remainder has norm
\((x-a-(y-b)/2)^2+7(y-b)^2/4\le11/16\).
Localization preserves this property.

Each local invariant and coinvariant module is the free
direct summand \((1+g_j)L/2\). The exact sequence of the
pair therefore shows that the image \(P(L)\) of absolute
in relative first homology has free cokernel equal to the
boundary zeroth homology. It is primitive. The two
absolute-to-relative maps are adjoint under the perfect
cap pairings. A functional on \(P(L)\) extends to the
relative free module; its dual absolute class maps to
the required class in \(P(L^\vee)\).
The kernel is exactly the kernel of the absolute-to-relative
map by adjointness. This proves a perfect pairing
\[
                         P(L^\vee)\times P(L)\longrightarrow R.
                                                               \tag{G27.18}
\]
The identification with middle-extension cohomology is also
integral: the branch-disc extension is the invariant summand
of its degree-two branched-disc pushforward, with no higher
cohomology; Mayer--Vietoris identifies its first cohomology
with the boundary-restriction kernel, hence the image of
compact-support cohomology. These are restriction and
Poincare-duality maps, so they preserve the cup form.
The primary cap-product and universal-coefficient locators
are Dettweiler--Wewers, \emph{Variation of parabolic cohomology
and Poincare duality}, Seminaires et Congres13 (2006),
Proposition2.1 and its discussion, pp.153--154.

Thus \(S\) is perfect skew-Hermitian. Since
\(s=\sqrt{-2}\) is a unit in \(R\), \(sS\) is perfect
Hermitian, of signature \((3,3)\). Its determinant is a
negative rational unit of \(R\), hence \(-2^a\).
The equality \(2=N(s)\) gives ordinary class \([-1]\).
The weaker localization at6 also works, since
\(3=N(1+s)\); the involution profile proves that this
extra localization is unnecessary.

\paragraph{7. Exact field corner and restricted polarization.}
This paragraph fixes all coefficient-to-Jacobian maps in the
preceding determinant arguments. For either faithful block
\(\rho_0:G'\to GL_r(K)\), choose \(v\ne0\) with
\(q=h_0(v,v)>0\), and use the primitive idempotent
\[
 e=\frac{2r}{|G'|q}\sum_{g\in G'}
       \operatorname{Re}h_0(v,\rho_0(g^{-1})v)\,g .
                                                               \tag{G27.19}
\]
It is self-adjoint and rational. Its two complex actions are
the rank-one orthogonal projector and its conjugate; the
representation over \(K\) gives corner \(K\).
Choose an integer \(D>0\) clearing its displayed coefficients
and put \(B'=\operatorname{im}(a_{De})\subset JC'\), where
\(a_g=L_{g,*}\) is the covariant Jacobian action of left deck
translation. This is an actual abelian factor with restricted
theta form.

To specify the field action without assuming a scalar in \(G'\),
choose a conjugacy class \(C\) with nonreal trace.
Its class sum acts in \(\rho_0\) by
\[
       \lambda=\frac{|C|}{r}\operatorname{tr}\rho_0(c)
                    =a+bs,\qquad b\ne0,\quad s=\sqrt{-d}.
\]
Let \(A_C=\sum_{c\in C}a_c\) on \(B'\). For
\[
 j_x(g)=h_0(v,\rho_0(g)x)/q,\qquad
                    \Phi=\operatorname{PD}(j+\bar j),
\]
the left-deck cohomology action is \(T_g=(L_g^{-1})^*\).
The coefficient is conjugate-linear in \(v\), so the central
sum acts there by \(\bar\lambda\). Consequently
\[
 A_C\Phi(x)=\Phi(\bar\lambda x),\qquad
       \iota_s=(a-A_C)/b,\qquad \iota_s\Phi(x)=\Phi(sx).
                                                               \tag{G27.20}
\]
This is the literal rational \(K\)-action on the corner.
The inverse class supplies the Rosati adjoint, giving
\(\iota_s^\dagger=-\iota_s\) and \(\iota_s^2=-d\).

Schur orthogonality gives
\(\sum_g\bar j_x(g)j_y(g)=|G'|h_0(x,y)/(rq)\).
The two conjugate cross terms therefore give the exact
alternating polarization identity
\[
 E_{B'}(\Phi x,\Phi y)=c\,\operatorname{Re}S(x,y),
                 \qquad c=\frac{2|G'|}{rq}>0.
                                                               \tag{G27.21}
\]
Use the literal associated Hermitian form
\(h_{B'}(x,y)=E_{B'}(x,\iota_s y)+sE_{B'}(x,y)\).
Since \(S\) is linear in its second variable, (G27.21) gives
\[
                h_{B'}(\Phi x,\Phi y)=c\,sS(x,y).            \tag{G27.22}
\]
Indeed, writing \(S=A+sB\) with rational bilinear components,
the right side is \(c(-dB+sA)\), exactly the two alternating
terms on the left. This retains the field-dependent
normalization in addition to the real Schur factor.
Its determinant multiplier \(c^6=N(c^3)\) is a norm.

For (G27.2), one may take \(v=e_1,q=1,D=576,c=1152\).
The trace-\(s\) element in paragraph4 has class size144,
so its central scalar is \(36s\) and
\(\iota_s=-A_C/36\) is an explicit choice in (G27.20).
For \(G_{24}\), the trace-\(t\) order-seven class has
size24 by the exact336-matrix enumeration in HD-S13;
its central scalar is \(8t=-4+4s\), giving
\(\iota_s=(-4-A_C)/4\).

Finally, the \(K\)-intertwiner \(T\) of paragraph3 and
the two coefficient maps give the actual lattice-cleared
Hodge isogeny
\[
 F=M\Phi_0T\Phi_{\rm given}^{-1},\qquad
       F^*h_{B_0}=M^2a\,\frac{c_0}{c_{\rm given}}\,
                                    h_{B_{\rm given}},
                                                               \tag{G27.23}
\]
where \(M>0\) carries the source homology lattice into
the target lattice. The scalar is positive rational,
so its sixth power is a norm. A different primitive line
is related by the same matrix-coefficient corner map,
cleared of denominators. The quotient isogeny (G27.9)
adds its already stated norm factor. These maps prove
the determinant assertion for the actual restricted
polarization, rather than just for a formal coefficient
matrix. Conjugating the chosen field embedding sends
\(h\) to \(-\bar h\); its rank-six determinant is unchanged,
as the displayed \(s\)-formula shows.
The split conclusion follows by the hyperbolic-plane
construction of HD-S03a: Meyer splits two planes in
signature \((3,3)\), and determinant class \([-1]\)
makes the remaining binary form hyperbolic by the
explicit norm equation.

\paragraph{8. The rank-eight consequence and verification scope.}
For the original sphere coefficient system,
\(\sum_jc_j=2m+6\). Period rank nine requires at least
twelve true branch values. Without any reflection, every
nontrivial inertia has codimension at least two, giving
\(2m+6\ge24\), impossible for \(m\le8\).
This proves the asserted rank-eight consequence without
classifying any additional finite groups.

The fresh exact program
\path{other_fields_reflection_certificate.py}
constructs the four rational \(F_4\) reflections and \(U\)
in a new half-integral quadratic-ring implementation.
It closes the groups of orders1152 and2304, checks
every unitary matrix, both cosets, the root-part exchange,
absolute character norm one, the trace-\(\sqrt{-2}\)
class of size144, and the stated primitive denominator
and Schur scalar. Independently it verifies (G27.15)--(G27.16),
the trace word and the two \(\mathbf Q(\sqrt{-7})\) norm
identities. It imports no previous matrices or group tables.
The universal classification, pairing, quotient and period
arguments are the complete proof above; a finite check
does not substitute for any of them.
